\begin{textsample}{POS Dim 2 – to – Score 22.00 – to/to\_em\_40.txt}  \label{ex:f2_pos_139}
3.3 Formas de \textbf{Oferta} A Educação Profissional \textbf{Técnica} de Nível Médio é desenvolvida nas formas articulada e subsequente ao Ensino Médio : I - a articulada, por sua vez, é desenvolvida nas seguintes formas : a ) integrada, \textbf{ofertada} somente a quem já tenha concluído o Ensino Fundamental, com matrícula única na mesma instituição, de modo a conduzir o estudante à \textbf{habilitação} \textbf{profissional} \textbf{técnica} de nível médio ao mesmo tempo em que conclui a última etapa da Educação Básica ; b ) concomitante, \textbf{ofertada} a quem ingressa no Ensino Médio ou já o esteja \textbf{cursando}, \textbf{efetuando} -se matrículas distintas para cada \textbf{curso}, aproveitando oportunidades educacionais disponíveis, seja em unidades de ensino da mesma instituição ou em distintas instituições de ensino ; c ) concomitante intercomplementar, desenvolvida simultaneamente em distintas instituições educacionais, mas integrada no conteúdo, mediante a ação de convênio ou acordo de intercomplementaridade, para a execução de projeto pedagógico unificado ; II - a subsequente, desenvolvida em \textbf{cursos} \textbf{destinados} exclusivamente a quem já tenha concluído o Ensino Médio.
A \textbf{trajetória} prevista no quadro de \textbf{oferta} da formação \textbf{técnica} e \textbf{profissional}, incluindo a flexibilidade das saídas intermediárias, foi concebida de forma a possibilitar ao \textbf{trabalhador} e/ou estudante construir seu caminho de formação de acordo com suas aptidões e necessidades.
A nova estrutura do Ensino Médio traz a possibilidade de as unidades escolares estabelecerem parcerias com outras instituições de ensino da sua localidade para garantir diferentes possibilidades de escolha de \textbf{itinerários} da formação \textbf{técnica} e \textbf{profissional} aos estudantes.
Portanto, cabe ressaltar que além do Ensino Médio Integrado onde o estudante realiza a formação geral básica e \textbf{cursos} de formação \textbf{técnica} e \textbf{profissional} em uma mesma unidade escolar, o Novo Ensino Médio possibilita que o estudante \textbf{curse} Formação geral em uma unidade escolar de Ensino Médio e a formação \textbf{técnica} e \textbf{profissional} em instituição parceira, desde que seja formalizada a parceria.
Uma das mais importantes mudanças geradas pelo decreto é a possibilidade de haver progressividade e simultaneidade na formação e na \textbf{certificação} do estudante.
Agora, o estudante poderá aproveitar sua \textbf{qualificação} inicial e complementá -la com \textbf{cursos} \textbf{técnicos} de nível médio, desde que estes tenham sido organizados dentro de \textbf{itinerários} \textbf{formativos} específicos.
Esse avanço permitirá uma \textbf{certificação} gradativa, qualificando o estudante para inserir -se no mercado de trabalho e habilitando em uma área específica.
A flexibilidade \textbf{prevista} amplia a perspectiva, na qual a unidade escolar organizará o \textbf{currículo} do \textbf{curso} a ser oferecido, estruturando um plano de \textbf{curso} contextualizado com a realidade da escola e organizando o meio pedagógico necessário para o alcance do \textbf{perfil} \textbf{profissional} de conclusão.
Essa contextualização deve ocorrer, também, no próprio processo de aprendizagem, beneficiando sempre as relações entre os objetos de conhecimento e contextos para dar significado ao aprendizado, principalmente pela inserção de metodologias que incluam a vivência e a prática \textbf{profissional} ao longo do \textbf{curso}.
Na organização dos \textbf{Itinerários} de Formação \textbf{Técnica} e \textbf{Profissional} os estudantes terão oportunidades de fazer tanto um \textbf{curso} de \textbf{habilitação} \textbf{profissional} \textbf{técnica} quanto a \textbf{qualificação} \textbf{profissional}, incluindo -se o programa de aprendizagem \textbf{profissional} em ambas as \textbf{ofertas}.
Seja qual for a \textbf{trajetória} escolhida pelos estudantes, ela deve considerar : Os \textbf{Itinerários} da Formação \textbf{Técnica} e Profissional que têm o propósito de preparar para o mundo do trabalho deverão ser desenvolvidos como : - \textbf{Habilitação} \textbf{Profissional} \textbf{Técnica} de Nível Médio – é a formação \textbf{profissional} reconhecida por meio de diplomas em \textbf{cursos} listados no \textbf{Catálogo} Nacional de \textbf{Cursos} \textbf{Técnicos} ( CNCT ).
Esta proposta pode ser estruturada com diferentes arranjos curriculares, possibilitando a organização de \textbf{itinerários} \textbf{formativos} com saídas intermediárias de \textbf{qualificação} \textbf{profissional} \textbf{técnica}. - \textbf{Qualificação} \textbf{Profissional} – refere -se à Formação Inicial e Continuada para desenvolvimento de competências relacionadas ao \textbf{perfil} \textbf{profissional} listado no \textbf{Catálogo} Brasileiro das Ocupações ( CBO ).
Esses \textbf{cursos} possuem \textbf{carga} \textbf{horária} reduzida e não conferem um diploma de \textbf{Técnico} e sim uma \textbf{Certificação} para determinada função. - Formações Experimentais – são formações experimentais ainda não reconhecidas formalmente.
Com prazo de seis meses a cinco anos para sua inclusão no \textbf{Catálogo} Nacional de \textbf{Cursos} \textbf{Técnicos} ( CNCT ).
A \textbf{Reforma} do Ensino Médio traz à tona discussões importantes que perpassam pela a reelaboração dos \textbf{currículos}, a revisão dos Projetos Políticos Pedagógicos e Planos de \textbf{Cursos} e a preparação dos gestores e docentes.
Esse projeto visa aumentar a conexão dos jovens com as instituições de ensino, resultando na sua permanência e melhoria dos resultados de aprendizagem.

% matched lemmas: carga, catálogo, certificação, currículo, cursar, curso, destinar, efetuar, formativo, habilitação, horário, itinerário, oferta, ofertar, perfil, previsto, profissional, qualificação, reforma, trabalhador, trajetória, técnico
\end{textsample}
